\begin{figure}[ht]
\centering
\begin{tikzpicture}[>=Stealth, font=\small, node distance=2.7cm]
  \node[draw, rounded corners, fill=blue!8, minimum width=3.4cm, minimum height=1cm] (x1) {$X_0(N)$};
  \node[draw, rounded corners, fill=green!8, right of=x1, minimum width=4.0cm, minimum height=1cm] (y) {$X_0(N;n)$};
  \node[draw, rounded corners, fill=blue!8, right of=y, minimum width=3.4cm, minimum height=1cm] (x2) {$X_0(N)$};

  \draw[->, thick] (y) -- node[above] {$\alpha_n$} (x1);
  \draw[->, thick] (y) -- node[above] {$\beta_n$} (x2);

  \node[draw, rounded corners, fill=orange!10, below=2.1cm of y, minimum width=6.0cm, minimum height=1cm] (j)
    {$J_0(N)$ 上的自同态 $T_n=(\beta_n)_*\alpha_n^*$};

  \draw[->, thick] ($(x1.south)+(0.1,0)$) -- ++(0,-0.85) -| (j.west);
  \draw[->, thick] ($(x2.south)+(-0.1,0)$) -- ++(0,-0.85) -| (j.east);

  \node[below=0.8cm of j] {在余切空间上，$H^0(J_0(N),\Omega^1)\cong \Sk_2(\GammaO(N))$，且 $T_n$ 的作用正是经典 Hecke 算子。};
\end{tikzpicture}
\caption{Hecke 算子来自模曲线之间的对应：先把数据拉回到 $X_0(N;n)$，再推出回 $X_0(N)$，于是得到 Jacobian 上的自同态。}
\label{fig:ch06-hecke-jacobian}
\end{figure}
